\textbf{Objective and budget controls.} At the high budget, the paired
Fed-minus-UNSIR RA difference is $0.063\pm0.371$ percentage points under IID
and $-0.463\pm0.361$ under Non-IID. At the low budget, the corresponding
differences are $-0.004\pm0.306$ and $-0.356\pm0.222$ points. Paired FA
differences are $0.200\pm1.735$ and $-0.233\pm0.451$ points at the high
budget, and $0.700\pm1.480$ and $-0.333\pm1.856$ at the low budget.
Thus, the common-budget controls do not support a consistent retained-utility
advantage of the unsquared whole-tensor penalty. In the three Non-IID pairs,
the squared-penalty adaptation has higher retained accuracy at both budgets.
The primary recipe's RA advantage therefore cannot be attributed to the
Fed-EMNU penalty alone. These descriptive results do not establish population
equivalence or isolate the separate effects of increasing noise and repair
budgets, which change together between low and high.

\textbf{Optimized versus Gaussian inputs.} Fed high minus Gaussian changes
FA by $-3.467\pm3.200$ points under IID, but by $2.500\pm12.322$ under
Non-IID. RA differences are $-0.430\pm0.107$ and $-0.330\pm0.195$ points.
Optimization improves mean class suppression in this IID control but does
not establish a robust benefit across partitions. The large Non-IID variation
and the absence of a retained-utility gain argue against claiming that noise
optimization is uniformly effective for these accuracy metrics. Target attack
AUC decreases relative to Gaussian inputs by $0.029\pm0.012$ under IID and
$0.022\pm0.018$ under Non-IID, with a decrease in all three pairs in each
partition. This is a favorable outcome for the declared attack, alongside the
mixed accuracy results; it does not establish a privacy guarantee. Neither noise choice demonstrates
verified deletion; the reported target attack and reference distances remain
necessary diagnostics.

\textbf{Event aggregation.} Equal retained counts make weighted and equal
event aggregation identical under IID; saved posterior arrays are exactly
equal in all three pairs. Under Non-IID, retained-count weighting increases
RA by $0.896\pm0.224$ points and FA by $0.600\pm0.557$ relative to equal
weighting. This supports a limited utility--suppression trade-off for the
stated event rule, not a general aggregation advantage. It does not resolve
the omission of separate repair contributions from non-target holders in
sparse-owner settings.

Taken together, the new controls strengthen attribution and expose negative
evidence. They support describing a particular empirical recipe and its
operating limits rather than a newly established superior noise principle.
Only three paired scientific seeds are used per partition; no multiple-testing
adjustment or formal equivalence claim is made.
